\section{Introducción a la teoría general de modelos lineales.}

\begin{prop}
    En un Modelo de Gauss-Markov de la forma $Y= X\beta + \varepsilon$, donde $Y$ es un vector de observaciones, $X$ es la matriz de diseño, $\beta$ es el vector de parámetros desconocidos y $\varepsilon$ es el vector de errores aleatorios, la estimación mínima cuadrática $\hat{\beta}$ del vector $\beta$ cumple:
    \begin{equation*}
        E\left[\hat{\beta}\right] = \beta, \quad Cov\left[\hat{\beta}\right] = \sigma^2 (X^TX)^{-1}.
    \end{equation*}
    \begin{proof}
        La estimación mínima cuadrática $\hat{\beta}$ se ha visto que es:
        \begin{equation*}
            \hat{\beta} = (X^TX)^{-1}X^TY.
        \end{equation*}

        Veamos entonces en primer lugar que es insesgada:
        \begin{align*}
            E\left[\hat{\beta}\right] &= E\left[(X^TX)^{-1}X^TY\right] = (X^TX)^{-1}X^TE\left[Y\right] 
            = (X^TX)^{-1}X^TE\left[X\beta + \varepsilon\right] \\
            &= (X^TX)^{-1}X^T\left[X\beta + \cancelto{0}{E\left[\varepsilon\right]}\right] = (X^TX)^{-1}X^TX\beta = \beta.
        \end{align*}

        Ahora veamos que su covarianza es $\sigma^2 (X^TX)^{-1}$:
        \begin{align*}
            Cov\left[\hat{\beta}\right] &= E\left[\left(\hat{\beta} - E\left[\hat{\beta}\right]\right)\left(\hat{\beta} - E\left[\hat{\beta}\right]\right)^T\right] \\
            &= E\left[\left(\hat{\beta} - \beta\right)\left(\hat{\beta} - \beta\right)^T\right]
        \end{align*}
        donde se ha aplicado que $E\left[\hat{\beta}\right] = \beta$. Calculemos cuánto vale $\hat{\beta} - \beta$:
        \begin{align*}
            \hat{\beta} - \beta &= (X^TX)^{-1}X^TY - \beta = (X^TX)^{-1}X^T\left(X\beta + \varepsilon\right) - \beta =\\&= (X^TX)^{-1}X^TX\beta + (X^TX)^{-1}X^T\varepsilon - \beta =\\
            &= \beta + (X^TX)^{-1}X^T\varepsilon - \beta = (X^TX)^{-1}X^T\varepsilon.
        \end{align*}

        Por tanto:
        \begin{align*}
            Cov\left[\hat{\beta}\right] &= E\left[\left((X^TX)^{-1}X^T\varepsilon\right)\left((X^TX)^{-1}X^T\varepsilon\right)^T\right] =\\
            &= E\left[(X^TX)^{-1}X^T\varepsilon \varepsilon^TX(X^TX)^{-1}\right] = (X^TX)^{-1}X^TE\left[\varepsilon \varepsilon^T\right]X(X^TX)^{-1} =\\
            &= (X^TX)^{-1}X^T\sigma^2 I_n X(X^TX)^{-1} = \sigma^2 (X^TX)^{-1}X^TX(X^TX)^{-1} = \sigma^2 (X^TX)^{-1}.
        \end{align*}
    \end{proof}
\end{prop}

\begin{prop}[Caracterización de las funciones estimables]
    Una función $\psi(\beta)$ es estimable si y solo si existe un vector $c \in \mathbb{R}^n$ tal que:
    \begin{equation*}
        \psi(\beta) = c^TX\beta.
    \end{equation*}
    \begin{proof}
        Por definición, una función $\psi(\beta)$ es estimable si existe un vector $c \in \mathbb{R}^n$ tal que:
        \begin{equation*}
            E\left[c^TY\right] = \psi(\beta).
        \end{equation*}

        Por tanto:
        \begin{align*}
            \psi(\beta) &= E\left[c^TY\right] = c^TE\left[Y\right] = c^TE\left[X\beta + \varepsilon\right] = c^T\left[X\beta + \cancelto{0}{E\left[\varepsilon\right]}\right] = c^TX\beta.
        \end{align*}
    \end{proof}
\end{prop}